\documentclass[10pt,a4paper]{amsart}
\usepackage[margin=19mm]{geometry}
\usepackage{amsmath,amssymb,mathtools}
\usepackage[scr=rsfs,cal=euler]{mathalfa}
\usepackage{xfrac}
\usepackage[unicode,hidelinks,bookmarks=false]{hyperref}
\newcommand{\ie}{i.e.\ }
\newcommand{\cf}{cf.\ }
\newcommand{\resp}{respectively\ }
\newcommand{\Subsection}{\subsection}
\providecommand{\nobreakdash}{\nobreak\hbox{-}}
\setlength{\emergencystretch}{2em}
\multlinegap0pt
\usepackage{amssymb}
\usepackage{bm}
\usepackage{caption}
\usepackage{xcolor}
\newtheorem{thm}{Theorem}
\newtheorem{lem}{Lemma}
\newtheorem{multicols}{Multicols}
\DeclareMathOperator{\E}{\mathbb{E}}
\DeclareMathOperator{\N}{\mathbb{N}}
\DeclareMathOperator{\bP}{\mathbb{P}}
\title{Scaling limits for some Mittag-Leffler queues}
\author{Giacomo Ascione, Luigia Caputo}
\date{Source: arXiv:2602.16521v1 (CC BY 4.0)}
\begin{document}
\maketitle

\noindent\emph{Source and licence.} Scaling limits for some Mittag-Leffler queues. Version arXiv:2602.16521v1 (CC BY 4.0), distributed by arXiv under the Creative Commons Attribution 4.0 International licence, http://creativecommons.org/licenses/by/4.0/, as recorded in the arXiv licence metadata for this exact version and shown on the abstract page https://arxiv.org/abs/2602.16521v1. Full licence notice: \texttt{licenses/arxiv-2602.16521-CC-BY-4.0.txt}.\par\smallskip

\noindent\textit{Editorial scope.} Counted unit: the article's lem lem:scalingMC2 (source0.tex lines 1717--1727). The article's complete original proof of it is delivered with it (source0.tex lines 1728--1776, one \texttt{proof} environment of 49 lines). No mathematics is written, repaired or re-derived: the statement and the proof are the article's own lines, in the article's own notation, and no retained line is substituted or re-flowed.  Retained uncounted as the source-local context the proof needs: lines 144--144; lines 225--232; lines 234--238; lines 240--242; lines 671--673; lines 675--688; lines 1713--1715; lines 2026--2033; lines 2047--2063; lines 2095--2100; lines 2182--2200.  Nothing was omitted from any retained span, so the package carries no omission record.  No retained line contains a citation, so the wrapper carries no bibliography and no bibliography entry.  No retained line contains a figure environment, an image-inclusion command or drawing code, so no image is delivered and the package declares no asset dependency.  Editorial formatting only, in addition: an AMS \texttt{amsart} preamble that restates the article's own declarations and macros used by the retained lines, namely the environments \texttt{thm}, \texttt{lem}, \texttt{multicols} on separate counters, together with the standard packages those lines need.  The retained environments print in this document as Theorem 1, Lemma 1, Multicols 1 in order of appearance, whereas the article prints the same items with the numbering of its own full document.  The complete native source of the article as distributed in the arXiv e-print for this exact version is bundled unchanged as \texttt{sources/194891/source0.tex}.
\subsection{Single server queueing systems}\label{subs:ssqs}

\begin{equation}\label{eq:mCren1}
	\bP(q_k=j \mid q_{k-1}=i)=\begin{cases}
		1 & j=1, \ i=0 \\
		p_{\lambda,\mu} & j=i+1 \not = 1 \\
		1-p_{\lambda,\mu} & j=i-1, \ i>0\\
		0 & |j-i|>1 \quad \mbox{ or } j=i \\
	\end{cases}
\end{equation}

\begin{equation}\label{eq:Y1}
	\bP(Y_k>t\mid q_{k-1})=\begin{cases} e^{-\lambda t} & q_{k-1}=0 \\
		e^{-(\lambda+\mu)t} & q_{k-1}>0.
	\end{cases}
\end{equation}

\begin{equation}\label{eq:countingproc}
	\widetilde{N}(t):=\max\left\{k \ge 0: \ \sum_{j=1}^{k}Y_k \le t\right\},
\end{equation}

\begin{equation}\label{eq:rescaling}
	\mathbf{Q}_n(t)=\frac{Q(nt)}{n^\delta}
\end{equation}

\begin{thm}\label{thm:classscal}
	Let $Q$ be the queue length process of a classical $M/M/1$ queue with inter-arrival and service rates $\lambda$ and $\mu$ and $\mathbf{Q}_n$ be its rescaling, defined in \eqref{eq:rescaling}. Let also $\rho=\frac{\lambda}{\mu}$ be the traffic intensity and $\mathbf{W}$ be a standard Brownian motion. Then:
	\begin{itemize}
		\item[$(i)$] If $\rho>1$ then for $\delta=1$ it holds $\mathbf{Q}_n \overset{\sf U}{\Rightarrow} (\lambda-\mu)\iota$;
		\item[$(ii)$] If $\rho=1$ then for $\delta=\frac{1}{2}$ it holds $\mathbf{Q}_n \overset{\sf U}{\Rightarrow} \sqrt{2\lambda}\Phi(\mathbf{W})$;
		%\item[$(iii)$] If $\rho<1$ then for $\delta=\frac{1}{2}$ it holds $\mathbf{Q}_n \overset{\sf U}{\Rightarrow} 0$.
		{\item[$(iii)$] If $\rho<1$ then for any $\delta>0$ it holds 
		%	$\mathbf{Q}_n \overset{\sf U}{\Rightarrow} 0$. In particular,
		\begin{equation*}
			\lim_{n \to +\infty}\sup_{t \ge 0}\mathbf{Q}_n(t)=0
			\end{equation*}
		almost surely.}
		\end{itemize}
\end{thm}

\begin{equation}\label{eq:scalingqn}
	\mathbf{q}_n(t)=\frac{q_{\lfloor nt \rfloor}}{n^\delta}.
\end{equation}

\begin{thm}\label{thm:contcomp}
	Let $(f_n)_{n \ge 1}\subset \mathbb{D}^k$ and $(g_n)_{n \ge 1}\subset \mathbb{D}_{\uparrow}$ with $(f_n,g_n) \overset{\sf M_1}{\to} (f,g) \in \mathbb{D}^{k+1}$. Assume that:
	\begin{itemize}
		\item[$(i)$] If $g(t) \in {\sf Disc}(f)$, then $t \not \in {\sf Disc}(g)$ and $g$ is strictly increasing in $t$;
		\item[$(ii)$] If $t \in {\sf Disc}(g)$, then $g(t),g(t-)\not \in {\sf Disc}(f)$ and the coordinates of $f$ are monotone in $[g(t-),g(t)]$. 
	\end{itemize}
	Then $f_n \circ g_n \overset{\sf M_1}{\to}f \circ g \in \mathbb{D}^k$.
\end{thm}

\begin{thm}\label{thm:contcounting}
	Let $(x_k^n)_{n,k \ge 1},(a_n)_{n \ge 1}\subset [0,+\infty)$ be such that $a_n>0$, $a_n \to +\infty$ and, defining the partial sums $(s_k^n)_{k \ge 0}$ as $s_0^n=0$ and $s_k^n=s_{k-1}^n+x_k^n$ for $k \ge 1$,
	%	\begin{equation*}
		%		s_0^n=0, \quad, s_k^n=\sum_{j=1}^{k}x_j^n,
		%	\end{equation*}
	it holds $s_k^n \to +\infty$ as $k \to +\infty$ for fixed $n \in \N$. Define $y_n(t)=s_{\lfloor t \rfloor}^n$ and $c_n$ the respective counting function. With $\mathbf{y}_n(t)=\frac{y_n(nt)}{a_n}$ and $\mathbf{c}_n(t)=\frac{c_n(a_nt)}{n}$, then the following statements are equivalent:
	\setlength{\columnsep}{3pt}
	\begin{multicols}{4}
		\begin{itemize}
			\item[$(i)$] $\mathbf{y}_n \overset{\sf M_1}{\to} \mathbf{y}$;
			\item[$(ii)$] $\mathbf{y}_n^{-1} \overset{\sf M_1}{\to} \mathbf{y}^{-1}$;
			\item[$(iii)$] $\mathbf{c}_n \overset{\sf M_1}{\to} \mathbf{y}^{-1}$;
			\item[$(iv)$] $\mathbf{c}_n^{-1}\overset{\sf M_1}{\to}\mathbf{y}$.
		\end{itemize}
	\end{multicols}
	\setlength{\columnsep}{10pt}
\end{thm}

\begin{thm}\label{thm:continuity}
	For all $f_1,f_2 \in \mathbb{D}$, for ${\sf T}={\sf U}$, ${\sf J_1}$ or ${\sf M_1}$ it holds
	\begin{equation*}
		d^{{\sf T}}(\Psi(f_1),\Psi(f_2)) \le d^{\sf T}(f_1,f_2), \qquad  d^{\sf T}(\Phi(f_1),\Phi(f_2)) \le 2d^{\sf T}(f_1,f_2).
	\end{equation*}
\end{thm}

\begin{thm}\label{thm:DTa}
	Let $\alpha \in (0,2]$. Assume that $X$ belongs to the normal domain of attraction of a stable law with parameters $C>0$ and $\beta \in [-1,1]$ if $\alpha<2$ or, if $\alpha=2$, set $C=\sqrt{{\sf Var}[X]}$. Let $(X_n)_{n \ge 1}$ be a sequence of i.i.d. random variables with the same distribution as $X$ and define the sequence of processes $(Y_n)_{n \ge 1}$ as  $Y_n(t)=\sum_{k=1}^{\lfloor nt \rfloor}X_k$. Then
	\begin{equation*}
		\left(\frac{C_\alpha}{C}\right)^{\frac{1}{\alpha}}n^{-\frac{1}{\alpha}}(Y_n-\iota_{m_n}) \overset{{\sf J}_1}{\Rightarrow} \mathbf{S},
	\end{equation*}
	where $\mathbf{S} \sim \mathbf{S}_\alpha(1,\beta)$ and
	\begin{equation}\label{eq:Ca}
		C_\alpha:=\begin{cases}
			\displaystyle \frac{1-\alpha}{\Gamma(2-\alpha)\cos\left(\frac{\pi \alpha}{2}\right)} & \displaystyle\substack{\alpha \in (0,2) \\ \alpha \not = 1} \\[7pt]
			\displaystyle 2/\pi & \alpha=1\\
			1 & \alpha=2
		\end{cases},
		\quad m_n=\begin{cases}
			0 & \alpha<1 \\[7pt]
			\displaystyle \frac{\pi Cn^{2}}{2}\E\left[\sin\left(\frac{2X}{\pi Cn}\right)\right] & \alpha=1 \\[7pt]
			n\E[X] & \alpha \in (1,2].
		\end{cases}
	\end{equation}
\end{thm}

\begin{lem}\label{lem:scalingMC2}
	Let $(q_k)_{k \ge 0}$ be a Markov chain as in \eqref{eq:mCren1}, $\mathbf{q}_n$ as in \eqref{eq:scalingqn} and $\mathbf{W}$ be a standard Brownian motion. Then
	\begin{itemize}
		\item[(i)] If $p_{\lambda,\mu}>\frac{1}{2}$ then for $\delta=1$ it holds $\displaystyle \mathbf{q}_n \overset{\sf U}{\Rightarrow}(2p_{\lambda,\mu}-1)\iota$
		\item[(ii)] If $p_{\lambda,\mu}=\frac{1}{2}$ then for $\delta=1/2$ it holds $\displaystyle \mathbf{q}_n \overset{\sf U}{\Rightarrow}\Phi(\mathbf{W})$
		\item[(iii)] If $p_{\lambda,\mu}<\frac{1}{2}$ then for any $\delta>0$ it holds 
		\begin{equation*}
		\lim_{n \to +\infty}\sup_{t \ge 0}\mathbf{q}_n(t)=0.
		\end{equation*}
	\end{itemize}
\end{lem}

\begin{proof}
	Let $(Y_k)_{k \ge 0}$ be as in \eqref{eq:Y1}. Let also $\widetilde{N}$ be the counting process of the $(Y_k)$, as in \eqref{eq:countingproc}. Then the process $Q(t)=q_{\widetilde{N}(t)}$ is the queue length process of a classical $M/M/1$ with traffic intensity $\rho=\frac{\lambda}{\mu}$. Let us immediately prove item (iii): since $\widetilde{N}:[0,+\infty) \to \N$ is surjective and $\rho<1$, we have
	\begin{equation*}
		\sup_{t \ge 0}\mathbf{q}_n(t)=n^{-\delta}\sup_{k \ge 0}q_k=n^{-\delta}\sup_{t \ge 0}Q(t)=\sup_{t \ge 0}\mathbf{Q}_n(t) \to 0.
	\end{equation*}
	To prove items (i) and (ii), with the notation of Section \ref{subs:ssqs}, let
	\begin{equation*}
		\mathbf{N}_n^{\sf a}(t)=n^{-1}N^{\sf a}(nt), \quad \mathbf{N}_n^{\sf d}(t)=n^{-1}N^{\sf d}(nt), \quad \mathbf{D}_n(t)=n^{-1}D(nt)
	\end{equation*}
	\begin{equation*}
		\mathbf{B}_n(t)=n^{-1}B(nt), \quad \mathbf{I}_n(t)=n^{-1}I(nt), \quad \mathbf{X}_n(t)=n^{-1}X(nt),
	\end{equation*}
	\begin{equation*}
	\mathbf{C}_n(t)=n^{-1}C(nt), \quad \mathbf{CS}_n(t)=n^{-1}CS_{\lfloor nt \rfloor}, \quad  \mathbf{G}_n(t)=n^{-1}G(t)
	\end{equation*}
	Then the following relations hold
	\begin{equation*}
		\mathbf{B}_n=\iota-\mathbf{I}_n, \quad \mathbf{I}_n=\Psi(\mathbf{X}_n), \quad \mathbf{X}_n=\mathbf{C}_n-\iota
	\end{equation*}	
	\begin{equation*}
		\mathbf{C}_n=\mathbf{CS}_n \circ \mathbf{N}_n^{\sf a}, \quad \mathbf{D}_n=\mathbf{N}^{\sf d}\circ \mathbf{B}_n, \quad \mathbf{G}_n=\mathbf{N}^{\sf a}_n+\mathbf{D}_n.
	\end{equation*}
	Once we recall the limit
	\begin{equation*}
		(\mathbf{N}_n^{\sf a}, \mathbf{N}_n^{\sf d}, \mathbf{CS}_n) \overset{\sf M_1}{\Rightarrow} (\lambda \iota, \mu \iota, \mu^{-1}\iota),
	\end{equation*}
	that clearly follows by Theorems \ref{thm:DTa} and \ref{thm:contcounting}, we immediately get
	\begin{equation*}
		\mathbf{C}_n \overset{\sf M_1}{\Rightarrow} \rho \iota, \quad \mathbf{X}_n \overset{\sf M_1}{\Rightarrow} (\rho-1) \iota.
	\end{equation*}
	Notice that we are only arguing for $\rho \ge 1$, hence $\rho-1 \ge 0$ and we have, by continuity of the regulator map as in Theorem \ref{thm:continuity},
	\begin{equation*}
		\mathbf{I}_n \overset{\sf M_1}{\Rightarrow} \mathbf{0}, \quad \mathbf{B}_n \overset{\sf M_1}{\Rightarrow} \iota, \quad \mathbf{D}_n \overset{\sf M_1}{\Rightarrow} \mu\iota, \quad 
		\mathbf{G}_n \overset{\sf M_1}{\Rightarrow} (\lambda+\mu) \iota.
	\end{equation*}
	Now let also
	\begin{equation*}
		\mathbf{J}_n(t)=n^{-1}J_{\lfloor nt \rfloor}
	\end{equation*}
	and use Theorem \ref{thm:contcounting} to guarantee that
	\begin{equation*}
		\mathbf{J}_n \overset{\sf M_1}{\Rightarrow} (\lambda+\mu)^{-1}\iota.
	\end{equation*}
	Now it remains to observe that $q_k=Q(J_k)$ for all $k \ge 0$, hence
	\begin{equation*}
		\mathbf{q}_n(t)=n^{-\delta}Q(J_{\lfloor nt \rfloor})=n^{-\delta}Q(n\mathbf{J}_n(t))=\mathbf{Q}_n \circ \mathbf{J}_n.
	\end{equation*}
	Since the limit of $\mathbf{Q}_n$ is continuous and the one of $\mathbf{J}_n$ is strictly increasing, we can use the continuous mapping theorem together with the fact that the composition is continuous in such a case (Theorem \ref{thm:contcomp}) to finally get the statement from Theorem~\ref{thm:classscal}.
\end{proof}

\end{document}
